
\section{Why I Am Leaving the Free Church}

“Folkebladet,” June 26, 1918

Quite simply: because I have come to realize that the manner in which I carry on congregational work is not in harmony with the view of congregational work maintained by the Guiding Principles of the Lutheran Free Church, but rather in harmony with the view of congregational work maintained by the large church bodies.

When the call came to take up the work in the newly established congregation in Detroit, Michigan, founded by Pastor P. M. Eriksen around New Year 1915, I accepted, though with considerable hesitation. The congregation had, after all, been founded by worldly people, which was contrary to the Guiding Principles of the Free Church, which I had learned to adhere to.

My hesitation about accepting the call, however, soon disappeared amid the busy work of clearing up this ecclesiastical wilderness. At the same time my method of work began to undergo serious changes. I began to see, 1) that although it was desirable to organize congregations in haste and hurry without further ado, a firm congregational organization was nevertheless an absolute necessity—especially in a large city, where people live so scattered and where there is so much that will lay claim to their interest—and that there must be no delay in bringing about as sound and vigorous a congregational organization as possible whenever outward circumstances favor such an organization. 2) that the matter here was to preach God’s Word purely and clearly, and then to make use of the materials that were at hand to build with.

The Lord has also shown us that he has taken up the work. Through the power of the Lord’s Word many have now been converted from their life in sin and live as Christians.

I became still more firmly established in my confidence in this method of work through the experience I gained in our mission work among our Scandinavian people in Toledo, Ohio (about 60 miles southwest of Detroit). I worked here for a year and a half without congregational organization. Several times the work was on the point of collapsing. And it was only through enormous exertions that it got under way again. Then we made the decision—we should have made it earlier. On November 11 last we organized a congregation with 50 adult members. Here also there are signs that the Lord will carry out his work through the working of the Word and the Holy Spirit.

In the light of what has been said above, it is clear that I have departed from the view of the proper course in church work maintained in the Guiding Principles of the Lutheran Free Church; and therefore, by continuing to stand in ecclesiastical fellowship with it, I do wrong both to the Free Church and to myself. At the same time, I have come nearer to sharing the view of the larger church body, and for that reason I have submitted my application for reception into the Norwegian Lutheran Church in America.

The congregation in Detroit was given an opportunity to decide for itself before I gave any indication of what I, for my own part, intended to do. The result of the vote was as follows: 98 voted to join the new church body, 4 against. Three of the four were, in truth, themselves inclined toward the larger church body; but they voted against the change because the Free Church was the branch of the church that had first begun the work there. The fourth voted against because he felt that things had been good as they were, and because he feared that the work of the congregation might suffer spiritually—a thing which I nevertheless hope and pray may not come to pass.

I wish here to express my thanks to the Free Church for all the good that has come to me through it. I leave it because I believe it is better and more advisable to stand wholly within the larger church body than to stand half within the Free Church. And yet, despite all this, I feel that, at bottom, I can never become anything other than a Free Church man—rightly understood.

L. O. Anderson.